\worldchapter{Konflikte}{combat}
\label{ch:konflikte}

\begin{multicols*}{2}

\section{Beginn des Konflikts}\label{sec:konflikt-beginn}\index{Konflikt}

Kommt es zu einem Konflikt, so verläuft die Zeit in Runden.
Zu Beginn des Konflikts wird eine grobe Karte gezeichnet, die eine Zonenstruktur enthält.
Distanzsegmente werden zwischen den Zonen und Gruppen festgelegt.
Dann wirft jeder Teilnehmer auf seine \textbf{Initiative} und alle Spielercharaktere auf ihre \textbf{Schnellreaktionen}.

Alle Teilnehmer beginnen den Konflikt mit \textbf{0 Aktionen}, es sei denn, ihnen gelingt die Probe auf eine Schnellreaktion.

\subsection{Initiative}\label{subsec:initiative}\index{Initiative}

Die Initiative bestimmt die Handlungsreihenfolge im Konflikt.
Hierzu wird eine Probe auf den Initiative-Zielwert abgelegt (\textbf{30} + \textbf{Geschick} + \textbf{Wahrnehmung}).
Die Teilnehmer handeln in der Reihenfolge der Punkte, die sie unter oder über ihrem Inititative-Wert geworfen haben.

\subsection{Schnellreaktion}\label{subsec:schnellreaktion}\index{Schnellreaktion}

Die Schnellreaktion ermöglicht es Spielercharakteren, eine Reaktion zu Beginn des Konflikts zu erhalten, um auf eventuelle Angriffe reagieren zu können.
Hierzu werfen die Spieler auf \textbf{Schnellreaktion}\index{Schnellreaktion}: \textbf{30} + \textbf{Instinkt} + \textbf{Wahrnehmung}.

Ist der Wurf erfolgreich, so erhält der Charakter eine Aktion zu Beginn des Konflikts, die ausschließlich als Reaktion eingesetzt werden kann.

\section{Aktionen pro Runde}\label{sec:aktionen-pro-runde}\index{Aktion}

Jeder Teilnehmer am Konflikt hat \textbf{2 Aktionen}\index{Aktion} pro Konfliktrunde.

\section{Rundenablauf}\label{sec:rundenablauf}

Jede Konfliktrunde hat ein Start- und ein Endfenster, in denen besondere Aktionen geschehen können.
Dazwischen handeln alle Charaktere in der Reihenfolge ihrer Initiative.

Kommt ein Charakter an die Reihe, so frischen sich seine Aktionen auf das Maximum auf.
Gehaltene Aktionen verfallen dabei.

\begin{enumerate}[leftmargin=*]
\item Startfenster der Runde
\item in Initiativreihenfolge: Aktionen auf Maximalwert auffrischen, dann ausgeben oder halten
\item Endfenster der Runde
\end{enumerate}

\section{Reaktionen}\label{sec:reaktionen}\index{Reaktion}

Im Konflikt ist es möglich, Aktionen zu halten, um später in der Kampfrunde zu handeln und auf Gegner zu reagieren.
Dies ist möglich, wenn der Charakter noch mindestens eine freie Aktion übrig hat und den Handelnden direkt wahrnehmen kann.

Für jeden Auslöser darf jeder Charakter einmal reagieren.
Es dürfen alle Charaktere reagieren, die die auslösende Handlung direkt wahrnehmen.
Reagieren mehrere Charaktere gleichzeitig, wird in Initiativreihenfolge aufgelöst.

Die Reaktion wird angesagt, sobald der Auslöser klar ist.
Sie wird vor der auslösenden Handlung abgehandelt.

\section{Aktionen}\label{sec:aktionsliste}\index{Aktion}

Aktionen können verwendet werden, um Handlungen auszuführen oder sich über Distanzen zu bewegen.
Die Bewegung innerhalb eines Distanzsegments benötigt keine Aktion.

Mögliche Aktionen in der eigenen Konfliktrunde sind:

\begin{description}[style=nextline,leftmargin=0pt]
\item[\textbf{Angreifen}] Nahkampfangriff
\item[\textbf{Schießen}] Fernkampfangriff
\item[\textbf{Bewegen}] 1 Distanzsegment bewegen
\item[\textbf{Sprinten}] bis zu 2 Distanzsegmente bewegen, danach -10 auf den nächsten Test der Runde
\item[\textbf{Deckung suchen}] bis zum Verlassen der Deckung +10 auf Reaktionen gegen Fernangriffe
\item[\textbf{Zielen}] der nächste Angriff erhält +10
\item[\textbf{Gegenstand benutzen}] Werkzeug, Verbrauchsgut, Nachladen, Mechanismus
\item[\textbf{Zaubern}] Zauber wirken, die keine Rituale sind
\item[\textbf{Unterstützen}] ein naher Verbündeter erhält +10 auf seine unmittelbare Probe
\item[\textbf{Zurückziehen}] Nahkampfbindung verlassen, ohne einen Gelegenheitsangriff auszulosen
\item[\textbf{Sich fassen}] Erschüttert entfernen oder einen sterbenden Verbündeten stabilisieren
\end{description}

\section{Bewegung im Konflikt}\label{sec:distanzsegmente}\index{Bewegung}\index{Distanzsegment}

Bewegung und Distanz wird durch \textbf{Distanzsegmente} dargestellt.
Ein Distanzsegment sind keine Meterangaben, sondern Entfernungskategorien zwischen Charakteren oder Gruppen.

Es gibt vier \textbf{Distanzsegmente}:

\begin{itemize}[leftmargin=*]
\item \textbf{Gebunden} – im Nahkampf mit einem Teilnehmer gebunden
\item \textbf{Nah} – innerhalb eines Raums oder einer begrenzten Fläche
\item \textbf{Fern} – im nächsten Raum, jenseits des Weges
\item \textbf{Weit entfernt} – weiter entfernt vom Geschehen
\end{itemize}

Bei Bewegung und der Bewertung von Distanzsegmenten gelten folgende Regeln:

\begin{itemize}[leftmargin=*]
\item \textbf{Bewegen} verschiebt 1 Distanzsegment.
\item \textbf{Sprinten} verschiebt bis zu 2 Distanzsegmente.
\item Der Wechsel von nicht \textbf{Gebunden} zu \textbf{Gebunden} bedeutet tatsächlichen Kontakt zu einem Gegner.
\item Umpositionierung im selben Distanzsegment ist frei, solange daraus kein deutlicher mechanischer Vorteil entsteht.
\item Falls Deckung, Flanke, Blocker oder Gefahr überwunden werden, erfordert dieselbe Umpositionierung \textbf{Bewegen} oder \textbf{Sprinten}.
\end{itemize}

Faustregeln für Innenräume:

\begin{itemize}[leftmargin=*]
\item in den Nachbarraum oder durch eine offene Verbindung wechseln: meist 1 Distanzsegment
\item in einen getrennten Raum, ein anderes Stockwerk oder hinter mehrere Sichtblocker wechseln: meist 2 Distanzsegmente
\end{itemize}

\section{Angriff abhandeln}\label{sec:angriff-auflosen}\index{Angriff}

Um einen Angriff durchzuführen, wirft der Angreifer eine passende Probe.

\begin{itemize}[leftmargin=*]
\item Nahkampf: \textbf{30} + \textbf{Körper} + \textbf{Kampf}
\item Fernkampf: \textbf{30} + \textbf{Geschick} + \textbf{Schießen}
\end{itemize}

Wenn der Verteidiger eine Aktion übrig hat, kann er mit folgenden Aktionen reagieren.

\begin{itemize}[leftmargin=*]
\item \textbf{Parieren}: \textbf{30} + \textbf{Instinkt} + \textbf{Kampf}
\item \textbf{Ausweichen}: \textbf{30} + \textbf{Geschick} + \textbf{Athletik}
\item \textbf{Blocken}: Schildbonus
\item \textbf{Treffer hinnehmen}: keine Reaktion
\end{itemize}

\subsection{Treffer und Schaden}\label{subsec:treffer-und-schaden}
\begin{itemize}[leftmargin=*]
\item Verliert der Angreifer die Gegenprobe: kein Treffer.
\item Nimmt der Verteidiger den Treffer und der Angriff misslingt: kein Treffer.
\item Nimmt der Verteidiger den Treffer und der Angriff gelingt: Grundschaden.
\item Gewinnt der Angreifer knapp: Grundschaden.
\item Gewinnt der Angreifer sauber oder stark: Grundschaden +1.
\item Ist der gewinnende Treffer kritisch: zusätzlich +1 Schaden.
\item Danach Rüstung einmal anwenden.
\end{itemize}

\section{Gelegenheitsangriff}\label{sec:parting-strike}\index{Gelegenheitsangriff}

Wenn jemand eine \textbf{Nahkampfbindung}\index{Nahkampfbindung} ohne \textbf{Zurückziehen} verlässt, darf ein gebundener Gegner sofort einen freien \textbf{Gelegenheitsangriff} ausführen.

Der Gelegenheitsangriff kostet keine Aktion, es ist eine freie Handlung.

\section{Zustandsverteilung im Konflikt}\label{sec:zustandsverteilung-im-kampf}

Waffen, Konsequenzen, Feuer, Gift oder spezielle Effekte können Zustände auslösen.
Das ist der Fall, wenn es explizit bei einer Waffe oder einem Effekt angegeben ist.
Ihre Regeln stehen gesammelt in \cref{ch:folgen}.

\end{multicols*}
